% 第 4 章 §4.3.3–§4.3.4 前段（印刷页 192–199）
% 承接：定理 4.3.5 的证明后半段（前段止于 ch04d 的式 (4.36)），故此处自「（续）」另起 proof 环境。
% 定理类计数器：本块首个环境为定义 4.3.9，故置 theorem = 8。
\setcounter{theorem}{8}

\begin{proof}[定理 4.3.5 的证明（续）]
令 $\mathcal{V}_\ell=\mathcal{V}$，并递归地定义 $\mathcal{V}_{i-1}$ 为 $\varphi_0(\mathcal{V}_i)$ 的饱和化。注意 $\mathcal{V}_{i-1}$ 是 $\mathcal{V}_i$ 的真子丛，否则 $\varphi_0$ 的限制将几乎处处是同构。最终该过程终止。调整 $\ell$ 使 $\mathcal{V}_0=\{0\}$、$\mathcal{V}_1\neq\{0\}$。由注 4.3.6，$\beta$ 保持滤过 $0=\mathcal{V}_0\subset\mathcal{V}_1\subset\cdots\subset\mathcal{V}_\ell=\mathcal{V}$。选取 $V$ 上的埃尔米特度量，并令 $\pi_i$ 为到 $V_i$ 的正交投影。注意
\[
I=\sum_{i=1}^{\ell}(\pi_i-\pi_{i-1})=\sum_{i=1}^{\ell}(\pi_i-\pi_i\pi_{i-1})=\sum_{i=1}^{\ell}\pi_i(I-\pi_{i-1}),
\]
且 $(I-\pi_i)\beta\pi_i=(I-\pi_{i-1})\varphi\pi_i=0$。于是
\[
\begin{aligned}
\operatorname{tr}(\beta\varphi_0)&=\operatorname{tr}(\varphi_0\beta)=\sum_{i=1}^{\ell}\operatorname{tr}\bigl(\varphi_0\beta\pi_i(I-\pi_{i-1})\bigr)\\
&=\sum_{i=1}^{\ell}\operatorname{tr}\bigl((I-\pi_{i-1})\varphi_0\beta\pi_i\bigr)=\sum_{i=1}^{\ell}\operatorname{tr}\bigl((I-\pi_{i-1})\varphi_0\pi_i\beta\pi_i\bigr)=0.
\end{aligned}
\]
故 $(\beta,\varphi_0)=0$，再由 (4.36) 推得 $[\beta]=0$。证明完成。
\end{proof}

\subsection{Corlette--Donaldson 定理}

\paragraph{埃尔米特度量与等变映射} 令 $D=\operatorname{SU}_n\backslash\operatorname{SL}_n(\C)$，$\rho\colon\pi\to\operatorname{SL}_n(\C)$。于是 $\pi$ 通过表示 $\rho$ 右作用在 $D$ 上。仿照 Donaldson，我们给出 $D$ 及其 $\operatorname{SL}_n(\C)$-作用的具体描述。令 $D$ 为所有满足 $\det M=1$ 的正定埃尔米特 $n\times n$ 矩阵 $M$ 构成的集合；此时右 $\operatorname{SL}_n$-作用由 $(M,g)\mapsto g^{-1}M(g^{-1})^*$ 给出。注意空间 $D$ 可解释为 $\C^n$ 上这样的埃尔米特内积所构成的空间：它在 $\det\C^n$ 上诱导出固定的内积。$D$ 上的不变度量由
\[
|M^{-1}dM|^2=\operatorname{tr}(M^{-1}dM)^2
\]
给出。

\begin{definition}
称映射 $u\colon\widetilde X\to D$ 为 \emph{$\rho$-等变}的，若对一切 $x\in X$、$\gamma\in\pi$ 都有 $u(x\gamma)=u(x)\rho(\gamma)$。
\end{definition}

令 $E=\widetilde X\times\C^n/\pi$。我们断言：$E$ 上的埃尔米特度量等价于 $\rho$-等变映射的选取，至多相差基点的选取。事实上，设 $u\colon\widetilde X\to D$ 是 $\rho$-等变的。按定义，$E$ 的截面是映射 $\sigma\colon\widetilde X\to\C^n$，满足 $\sigma(x\gamma)=\sigma(x)\rho(\gamma)$。于是，若定义 $\|\sigma\|^2(x)=\langle\sigma(x),\sigma(x)u(x)\rangle_{\C^n}$，则
\[
\|\sigma\|^2(x\gamma)=\langle\sigma(x)\rho(\gamma),\sigma(x)u(x)(\rho(\gamma)^{-1})^*\rangle_{\C^n}=\|\sigma\|^2(x),
\]
因而这是 $E$ 上定义良好的度量。反之，给定度量 $H$，若 $\sigma_i$ 是截面，则对某个埃尔米特矩阵值函数 $u(x)$ 记 $\langle\sigma_i,\sigma_j\rangle_H(x)=\langle\sigma_i(x),\sigma_j(x)u(x)\rangle_{\C^n}$。于是
\[
\begin{aligned}
\langle\sigma_i(x),\sigma_j(x)u(x)\rangle_{\C^n}&=\langle\sigma_i,\sigma_j\rangle_H(x)=\langle\sigma_i,\sigma_j\rangle_H(x\gamma)\\
&=\langle\sigma_i(x)\rho(\gamma),\sigma_j(x)\rho(\gamma)u(x\gamma)\rangle_{\C^n}\\
&=\langle\sigma_i(x),\sigma_j(x)\rho(\gamma)u(x\gamma)\rho(\gamma)^*\rangle_{\C^n}
\end{aligned}
\]
对一切截面成立。因此 $\rho(\gamma)u(x\gamma)\rho(\gamma)^*=u(x)$，即 $u$ 是 $\rho$-等变的。

\paragraph{调和度量} 若 $u\colon\widetilde X\to D$ 是连续可微的 $\rho$-等变映射，我们如下定义其\emph{能量}（energy）。导数 $du$ 是 $T^*\widetilde X\otimes u^*(TD)$ 的截面。我们在 $D$ 上固定了不变度量，故范数为 $e_u(x)=|du|^2(x)$。事实上，由等变性，$e_u(x)$ 在 $\pi$ 下不变，因而给出 $X$ 上定义良好的函数，称为\emph{能量密度}（energy density）。$u$ 的能量按定义即为
\begin{equation}
E_\rho(u)=\int_X e_u(x)\,\omega.
\label{eq:4.37}\tag{4.37}
\end{equation}
注意能量只依赖于 $X$ 上的共形结构，而与完整度量无关。

$E_\rho$ 的 Euler--Lagrange 方程容易写出。定义
\begin{equation}
\tau(u)=d^*_{\nabla}du.
\label{eq:4.38}\tag{4.38}
\end{equation}
上式中我们注意到，丛 $u^*(TD)$ 带有联络 $\nabla$：即 $D$ 上 Levi--Civita 联络的拉回。$d_{\nabla}$ 正是相对于此联络定义的。张量 $\tau(u)$ 称为\emph{张力场}（tension field），它是 $u^*(TD)$ 的截面。

\begin{definition}
称 $C^2$ 的 $\rho$-等变映射 $u$ 为\emph{调和}（harmonic）的，若它满足
\begin{equation}
\tau(u)=0.
\label{eq:4.39}\tag{4.39}
\end{equation}
\end{definition}

方程 (4.39) 是关于 $u$ 的二阶椭圆非线性偏微分方程。这一说法略有误导，因为 $u$ 是映射，而不是一组函数。这一令人不快的事实使弱解（weak solution）的定义稍显棘手。在紧流形之间的映射的情形（非等变问题），规避此困难的一种办法是取靶空间到欧氏空间的 Nash 等距嵌入（Nash isometric embedding），再用坐标函数改写方程（参看 [53]）。适合等变问题的另一种更精细的技巧，是内蕴地定义 Sobolev 空间理论（参看 [41, 45, 46]）。另一方面，若假定 $u$ 是 Lipschitz 连续的，则可引入 $D$ 上的局部坐标 $\{y^a\}$ 并局部地写出 (4.39)。由 Rademacher 定理（Rademacher's theorem），拉回 $s_a=u^*(\partial/\partial y^a)$ 几乎处处给出 $u^*(TD)$ 的局部标架，而 $\nabla$ 在此标架下的联络形式为 $\Gamma^c_{ab}(u)du^a\otimes s_c$，其中 $\Gamma^c_{ab}(u)$ 是 $D$ 上的 Christoffel 符号沿 $u$ 的取值。把 $u=(u^1,\dots,u^N)$ 按 $\{y^a\}$ 上的坐标写出，容易看出 (4.39) 的局部表达式成为
\begin{equation}
-\tau(u)^a=\Delta u^a+\Gamma^a_{bc}(u)\nabla u^b\cdot\nabla u^c=0.
\label{eq:4.40}\tag{4.40}
\end{equation}
为明确起见，第二项中的点乘指 $X$ 上的度量，而 $\Delta$ 是 $X$ 上的 Laplace 算子。注意此方程关于 $X$ 上的度量是共形不变（conformally invariant）的，这正是能量泛函（energy functional）本身共形不变这一事实的体现。

由上一节，$\rho$-等变映射等价于埃尔米特度量的选取。给定平坦联络 $\nabla$ 与 $E$ 上的埃尔米特度量，我们可以用更内蕴的方式构造等变映射。首先把 $\nabla$ 与 $E$ 提升，得到万有覆叠 $\widetilde X$ 上平凡丛上的平坦联络。我们用同一记号表示提升后的丛与联络。若选取覆盖基点 $p$ 的基点 $\hat p$（对应 $\pi_1(X,p)$），并为纤维 $E_{\hat p}$ 选取酉标架 $\{e_i(\hat p)\}$，令 $\{e_i(x)\}$ 由关于 $\nabla$ 的平行移动给出。则映射 $u\colon\widetilde X\to D$ 由 $x\mapsto\langle e_i,e_j\rangle(x)$ 给出。它是 $\rho$-等变的，且除 $\hat p$ 与 $D$ 中基点的选取之外唯一确定。

反之，若 $u\colon\widetilde X\to D$ 是任一满足 $u(\hat p)=I$ 的 $\rho$-等变映射，则它定义了埃尔米特度量，而 $u$ 正是如上构造的该度量的等变映射。

注意存在我们已在 Higgs 丛中见过的等价性。若 $g\in \mathcal{G}^{\C}_E(p)$，则由对 $(g(\nabla),H)$ 得到的 $\rho$-等变映射与由 $(\nabla,Hg)$ 得到的相同。最后，若用常数 $g\in\operatorname{SL}_n(\C)$ 作用，同样成立，只是映射此时变为 $(\rho\cdot g)$-等变的。此事的要旨是：求调和度量等价于在 $\nabla$ 的 $\mathcal{G}^{\C}_E$ 轨道内求调和等变映射。

给定数据 $(\nabla,H)$，我们可唯一地写 $\nabla=d_A+\Psi$，其中 $d_A$ 是 $(E,H)$ 上的酉联络，$\Psi$ 是取值于埃尔米特自同态丛 $\sqrt{-1}\mathfrak{g}_E$ 中的 $1$-形式。相对于局部标架 $\{s_i\}$，可将 $\Psi$ 显式定义为
\begin{equation}
\langle\Psi s_i,s_j\rangle=\frac{1}{2}\bigl\{\langle\nabla s_i,s_j\rangle+\langle s_i,\nabla s_j\rangle-d\langle s_i,s_j\rangle\bigr\}.
\label{eq:4.41}\tag{4.41}
\end{equation}

\begin{lemma}[{参看 [20]}]
上述映射的能量由 $E_\rho(u)=4\|\Psi\|^2$ 给出。
\end{lemma}

\begin{proof}
由上述定义以及 $d_A$ 是酉的，可得 $du_{i,j}=\langle d_Ae_i,e_j\rangle+\langle e_i,d_Ae_j\rangle$。另一方面，$e_i$ 关于 $\nabla$ 平行，故 $d_Ae_j=-\Psi e_j$。于是 $u^{-1}du=-2\Psi$。
\end{proof}

\begin{definition}
称 $H$ 为\emph{调和度量}，若上述映射 $u$ 是调和映射。
\end{definition}

\begin{proposition}[{Corlette, [11]}]
若 $\rho$ 容许调和度量，则 $\rho$ 是半单（semisimple）的。
\end{proposition}

\begin{proof}
设 $H$ 是临界度量，但 $\nabla$ 可约。令 $V_1\subset V$ 是关于联络 $\nabla$ 不变的子丛，$V_2$ 为 $V_1$ 的正交补，$H_1$、$H_2$ 为诱导度量。我们可以写
\[
\nabla=\begin{pmatrix}\nabla_1&\beta\\0&\nabla_2\end{pmatrix}=\begin{pmatrix}d_{A_1}+\Psi_1&\beta\\0&d_{A_2}+\Psi_2\end{pmatrix},
\]
其中 $\beta\in\Omega^1(X,\operatorname{Hom}(V_2,V_1))$。只需证明联络分裂，换言之 $\beta\equiv0$；命题随即由归纳法得到。现在由 (4.41) 可知，若 $s_1,s_2$ 是 $V_1$ 的局部截面，则 $\langle\Psi s_1,s_2\rangle=\langle\Psi_1s_1,s_2\rangle$。类似地，对 $V_2$ 的局部截面有 $\langle\Psi s_1,s_2\rangle=\langle\Psi_1s_1,s_2\rangle$。\begin{annote}原文此处亦作 $\langle\Psi_1s_1,s_2\rangle$，与上一句完全相同；但按下面 $\Psi$ 的分块形状，对 $V_2$ 的局部截面应为 $\langle\Psi_2s_1,s_2\rangle$。译文照排原文。[来源：上下文推断]\end{annote}另一方面，若 $s_i\in V_i$，则 $\langle\Psi s_1,s_2\rangle=\frac{1}{2}\langle s_1,\beta s_2\rangle$。由此可得
\[
\Psi=\begin{pmatrix}\Psi_1&\frac{1}{2}\beta\\[2pt]\frac{1}{2}\beta^*&\Psi_2\end{pmatrix}.
\]
我们现在把度量 $H$ 形变为族 $H_t$ 如下：缩放 $H_1\mapsto e^{-(\operatorname{rank}V_2)t}H_1$，$H_2\mapsto e^{+(\operatorname{rank}V_1)t}H_2$。这当然保持正交分裂与条件 $\det H_t=1$。但 $H_t$ 是 $\rho$-等变映射的测地同伦（geodesic homotopy），故由 Hartman 的一个结果，能量 $E_\rho(u_t)$ 是凸的 [32]。另一方面，由引理 4.3.11，
\[
\frac{1}{4}E_\rho(u_t)=\|\Psi_1\|^2_{H_1}+\|\Psi_2\|^2_{H_2}+\|\beta\|^2_H\,e^{-(\operatorname{rank}V)t/2}.
\]
特别地，$E_\rho(u_t)$ 当 $t\to\infty$ 时保持有界。$E_\rho(u_t)$ 在 $t=0$ 处有临界点的唯一可能是 $E_\rho(u_t)$ 为常数，而这蕴涵 $\beta\equiv0$。证明完成。
\end{proof}

\paragraph{Corlette--Donaldson 定理} 本节我们证明如下定理。

\begin{theorem}[{Corlette [11]，Donaldson [20]，Jost--Yau [42]，Labourie [47]}]
设 $\rho\colon\pi\to\operatorname{SL}_n(\C)$ 半单。则存在 $\rho$-等变的调和映射 $u\colon\widetilde X\to D$。
\end{theorem}

下述结果可与引理 4.2.26 相比较。它将在下文讨论调和映射流时证明。

\begin{lemma}
对任意 $\rho\colon\pi\to\operatorname{SL}_n(\C)$，存在 $\rho$-等变映射的序列 $u_j\colon\widetilde X\to D$，满足下列条件：
\begin{enumerate}[label=(\roman*),leftmargin=3.2em,itemsep=1pt]
\item $u_j$ 是能量极小（energy minimizing）的；
\item $u_j$ 具有一致有界的 Lipschitz 常数；
\item $\tau(u_j)\to0$ 于 $L^2$ 中。
\end{enumerate}
\end{lemma}

\begin{lemma}
设 $\rho\colon\pi\to\operatorname{SL}_n(\C)$ 不可约，并设 $u_j\colon\widetilde X\to D$ 是具有一致 Lipschitz 常数的 $\rho$-等变映射序列。则 $u_j(\hat p)$ 有界。
\end{lemma}

\begin{proof}
若不然。令 $h_j=u_j(\hat p)$，并选取 $\varepsilon_j\to0$，使得（或许在取子列之后）$\varepsilon_jh_j\to h_\infty\neq0$。注意 $\det h_\infty=0$，故 $V=\ker h_\infty$ 是 $\C^n$ 的真子空间。我们断言 $\rho(\pi)$ 固定 $V$。事实上，若 $\rho(\gamma)=g^{-1}$ 且 $v\in V$，则由 $d(u_j(\hat p),u_j(\hat p)\cdot g^{-1})$ 一致有界，有
\[
|\langle w,vh_j\rangle_{\C^n}-\langle w,vgh_jg^*\rangle_{\C^n}|\le B,
\]
其中 $B$ 是与 $j$ 无关的常数，且对一切 $w\in\C^n$ 成立。由此推出
\[
|\langle w,v\varepsilon_jh_j\rangle_{\C^n}-\langle wg,vg\varepsilon_jh_j\rangle_{\C^n}|\longrightarrow0
\]
又因 $vh_\infty=0$，我们得到 $\langle wg,vgh_\infty\rangle_{\C^n}=0$。由于 $w$ 任意，故 $vg\in V$。
\end{proof}

\begin{proof}[定理 4.3.14 的证明]
由归纳法，只需对不可约表示证明结论。设 $u_j$ 是如引理 4.3.16 中的极小化序列，其存在性由引理 4.3.15 保证。由 Ascoli 定理（Ascoli's theorem），存在一致收敛的子列，仍记作 $u_j$，其极限 $u_\infty$ 是 Lipschitz 的 $\rho$-等变映射。我们断言可安排 $u_\infty$ 为调和映射。事实上，由于收敛是一致的，可取局部坐标并写出 $u^a$。于是，由于 $|du^a|$ 一致有界，可进一步假定 $u_j$ 在 $L^2_{1,\mathrm{loc}}$ 中弱收敛到 $u_\infty$。由引理 4.3.15 (iii) 的条件，坐标 $u^a_\infty$ 属于 $L^2_{1,\mathrm{loc}}$，并构成 (4.40) 的弱解。由于 $u_\infty$ 是 Lipschitz 的，Laplace 算子的椭圆正则性蕴涵 $u_\infty\in L^2_{2,\mathrm{loc}}$。由 (4.40) 之后的附注，可假定 $X$ 上的局部度量为欧氏的。现在求导得到
\[
\begin{aligned}
\Delta(\nabla u^a_\infty)+\nabla\bigl(\Gamma^a_{bc}(u_\infty)\nabla u^b_\infty\cdot\nabla u^c_\infty\bigr)&=0,\\
\Delta(\nabla^2u^a_\infty)+\nabla^2\bigl(\Gamma^a_{bc}(u_\infty)\nabla u^b_\infty\cdot\nabla u^c_\infty\bigr)&=0.
\end{aligned}
\]
注意，由于 $u_\infty$ 是 Lipschitz 的，第一式中的第二项属于 $L^2$。于是 $u^a_\infty\in L^2_{3,\mathrm{loc}}$。由于包含 $L^2_3\hookrightarrow L^4_2$，上式中第二式的第二项也属于 $L^2$。这又蕴涵 $u^a_\infty\in L^2_{4,\mathrm{loc}}$。最后，$L^2_4\subset C^{2,\alpha}$，故 $u_\infty$ 是调和映射方程 (4.39) 的强解（strong solution）。证明完成。
\end{proof}

\paragraph{调和映射流} 调和映射流定义为
\begin{equation}
\dot u=-\tau(u).
\label{eq:4.42}\tag{4.42}
\end{equation}
这里 $u_t$ 是一族 $\rho$-等变映射。由于 $D$ 具有非正曲率，该流的性态很好。长时间存在性在 [21, 30] 中证明。

能量沿流的变分由
\[
\frac{d}{dt}E(u_t)=2\int_X\langle du,d\dot u\rangle=2\int_X\langle d^*_\nabla du,\dot u\rangle\omega=-2\int_X|\tau(u)|^2\omega.
\]
给出。特别地，能量沿流递减。此外，
\begin{equation}
2\int_0^\infty dt\int_X|\tau(u_t)|^2\omega\le E(u_0).
\label{eq:4.43}\tag{4.43}
\end{equation}

我们现在可以证明引理 4.3.15 了。

\begin{proof}[引理 4.3.15 的证明]
证明基于关于能量密度沿调和映射流变化的有名 Eells--Sampson--Bochner 公式 [21]。设 $u=u(t,x)$ 是 (4.42) 的解，$e=e_u(t,x)$。则
\[
-\frac{\partial e}{\partial t}+\Delta e=|\nabla du|^2+\operatorname{Ric}_X(du,du)-\operatorname{Riem}_D(du,du,du,du).
\]
现在由于 $\operatorname{Riem}_D\le0$ 且 $\operatorname{Ric}_X$ 有负常数下界，有
\[
\frac{\partial e}{\partial t}-\Delta e\le C\cdot e.
\]
用热核（heat kernel）作显式计算，此不等式连同能量递减的事实蕴涵如下类型的估计，例如对 $t\ge1$，
\begin{equation}
\sup e_{u_t}\le C\cdot E_{u_0},
\label{eq:4.44}\tag{4.44}
\end{equation}
其中 $C$ 只依赖于 $X$ 与 $D$ 的几何。

现在令 $u^{(j)}$ 是 $\rho$-等变映射的能量极小序列。令 $u^{(j)}_t$ 为经 (4.42) 的时间 $t$ 流之后对应的映射。于是，由于能量沿流递减，对任意序列 $t_j$ 的选取，$u^{(j)}_{t_j}$ 也是能量极小的。另一方面，(4.44) 的右端一致有界，故若取每个 $t_j\ge1$，则 $u^{(j)}_{t_j}$ 也一致 Lipschitz。最后，对每个固定的初始条件 $u_0$，(4.43) 蕴涵沿某序列 $\tau(u_{t_j})\to0$ 于 $L^2$ 中。由对角线论证（diagonalization argument）可使 $u^{(j)}_{t_j}$ 也满足此性质。
\end{proof}

\subsection{超凯勒约化}

\paragraph{模空间彼此实同构} 利用 (4.41) 并给定埃尔米特度量，我们可以把全体联络的空间等同于
\[
\mathcal{C}_E=\bigl\{(A,\Psi)\in\mathcal{A}_E\times\Omega^1(M,\sqrt{-1}\mathfrak{g}_E)\bigr\}.
\]\begin{annote}原文此处作 $\Omega^1(M,\sqrt{-1}\mathfrak{g}_E)$，而全章其余各处均以 $X$ 记该 Riemann 面，$M$ 未作定义；译文照排原文。[来源：上下文推断]\end{annote}
于是 $\mathcal{C}_E$ 是\emph{超凯勒流形}（hyperkähler manifold），而规范群 $\mathcal{G}_E$ 的作用带有相伴的矩映射
\begin{equation}
\mu_1(A,\Psi)=F_A+\frac{1}{2}[\Psi,\Psi],\qquad \mu_2(A,\Psi)=d_A\Psi,\qquad \mu_3(A,\Psi)=d_A(*\Psi).
\label{eq:4.45}\tag{4.45}
\end{equation}\begin{annote}原文此处作 “the gauge group $\mathcal{G}$”，与本页后文及全节的 $\mathcal{G}_E$、$\mathcal{G}^{\C}_E$ 记号不一致；译文按全章记号统一作 $\mathcal{G}_E$。[来源：上下文推断]\end{annote}
令 $\mathbf{m}=(\mu_1,\mu_2,\mu_3)$。超凯勒商按定义为
\[
\mathbf{m}^{-1}(0)\big/\mathcal{G}_E=\mu_1^{-1}(0)\cap\mu_2^{-1}(0)\cap\mu_3^{-1}(0)\big/\mathcal{G}_E.
\]
我们上面讨论的两个图景相当于分步约化 $\mathcal{C}_E$，但方式有两种。第一种是 Hitchin 与 Simpson 的观点，见 §4.2.3。即，Higgs 丛的空间由
\[
\mathcal{B}_E=\mu_2^{-1}(0)\cap\mu_3^{-1}(0)\subset\mathcal{C}_E
\]
给出，其中 $\Psi$ 与 $\Phi$ 的关系是 $\Psi=\Phi+\Phi^*$，反之 $\Phi$ 是 $\Psi$ 的 $(1,0)$ 部分。正如曲面上的函数的情形，$\Psi$ 调和当且仅当 $\Phi$ 全纯。现在定理 4.2.17 保证每个多稳定 Higgs 丛的轨道都与 $\mathcal{B}^{ss}$ 内的零点集 $\mu_1^{-1}(0)$ 相交。因此有
\[
\mathfrak{M}^{(n)}_D=\mathcal{B}^{ss}_E\big/\!\!\big/\mathcal{G}^{\C}_E=\mathbf{m}^{-1}(0)\big/\mathcal{G}_E=\mu_1^{-1}(0)\cap\mu_2^{-1}(0)\cap\mu_3^{-1}(0)\big/\mathcal{G}_E.
\]
第二种观点（例如 Corlette 与 Donaldson，见 §4.3.3）来自如下观察：平坦联络的空间是
\[
\mathcal{C}^{\mathrm{flat}}_E=\mu_1^{-1}(0)\cap\mu_2^{-1}(0)\subset\mathcal{C}_E.
\]
给定 $\nabla\in\mathcal{C}^{\mathrm{flat}}_E$，相伴的 $\mathrm{hol}(\nabla)$-等变映射为调和这一条件恰为 $\nabla\in\mu_3^{-1}(0)$。事实上，设 $\delta\nabla$ 是 $\nabla$ 的变分。由 (4.41) 推出 $\delta\Psi=\delta\nabla+(\delta\nabla)^*$。在复规范变换 $g^{-1}\delta g=\varphi$ 的情形，有 $\delta\nabla=\nabla\varphi$，且 $\delta\Psi=d_A(\varphi+\varphi^*)+[\Psi,\varphi-\varphi^*]$。容易看出第二项在变分 $\operatorname{tr}(\delta\Psi\wedge*\Psi)+\operatorname{tr}(\Psi\wedge*\delta\Psi)$ 中不贡献（直接计算可知，也可由酉规范变换不改变相伴等变映射这一事实看出）。所以，由引理 4.3.11 有
\[
\begin{aligned}
\delta E(u)&=4\int_X\operatorname{tr}(\delta\Psi\wedge*\Psi)+\operatorname{tr}(\Psi\wedge*\delta\Psi)\\
&=4\int_X\operatorname{tr}\bigl(d_A(\varphi+\varphi^*)\wedge*\Psi\bigr)+\operatorname{tr}\bigl(\Psi\wedge*d_A(\varphi+\varphi^*)\bigr)\\
&=-8\int_X\operatorname{tr}\bigl((\varphi+\varphi^*)d_A(*\Psi)\bigr).
\end{aligned}
\]
由于 $\Psi$ 是埃尔米特的而 $\varphi$ 任意，$\Psi$ 是能量的临界点当且仅当 $d_A(*\Psi)=0$。

现在定理 4.3.14 保证每个半单表示的轨道包含调和映射。因此和乐映射（holonomy map）给出同胚
\[
\mathfrak{M}^{(n)}_B\simeq\mathcal{C}^{\mathrm{flat}}_E\big/\!\!\big/\mathcal{G}^{\C}_E\simeq\mu_1^{-1}(0)\cap\mu_2^{-1}(0)\cap\mu_3^{-1}(0)\big/\mathcal{G}_E.
\]
所以，Dolbeault 模空间与 Betti 模空间确实一致！

\begin{theorem}[{Simpson，[60, 61]}]
上述描述的对应给出同胚 $\mathfrak{M}^{(n)}_D\simeq\mathfrak{M}^{(n)}_B$。
\end{theorem}

\paragraph{等变上同调} 与 YMH（Yang--Mills--Higgs）流的情形一样，调和映射流在 $t\to\infty$ 时实际上具有连续性。为描述这一点，令 $\mathcal{G}_E(p)\subset \mathcal{G}_E$ 记在点 $p$ 处为恒同的规范变换所构成的子群。现在和乐映射给出紧嵌入（proper embedding）
\begin{equation}
\mathrm{hol}\colon\mathbf{m}^{-1}(0)/\mathcal{G}_E(p)\hookrightarrow\operatorname{Hom}(\pi,\operatorname{SL}_n(\C)),
\label{eq:4.46}\tag{4.46}
\end{equation}
它是 $\operatorname{SU}_n$-等变的。

\begin{theorem}[{参看 [15]}]
包含 (4.46) 是 $\operatorname{SU}_n$-等变的形变收缩。
\end{theorem}

显式的收缩是这样给出的：用调和映射流在表示空间上定义一个流。固定 $p$ 的提升 $\tilde p\in\widetilde X$。给定 $\rho\in\operatorname{Hom}(\pi,\operatorname{SL}_n(\C))$，选取 $\nabla\in\mathcal{C}^{\mathrm{flat}}_E$ 使 $\mathrm{hol}(\nabla)=\rho$。埃尔米特度量给出唯一的 $\rho$-等变提升 $u\colon\widetilde X\to D$，满足 $u(\tilde p)=I$。令 $u_t$、$t\ge0$ 为以 $u$ 为初始条件的 (4.42) 的解。存在唯一的连续族 $h_t\in\operatorname{SL}_n(\C)$、$h_t^*=h_t$，满足 $h_0=I$ 且 $h_tu_t(\tilde p)=z$。\begin{annote}原文未定义 $z$；按上下文，它应指 $D$ 中某个固定基点，$h_t$ 的作用是把 $u_t(\tilde p)$ 拉回该点。译文照排原文。[来源：上下文推断]\end{annote}注意，若选取另一个满足 $\mathrm{hol}(\tilde\nabla)=\rho$ 的平坦联络 $\tilde\nabla$，则它与 $\nabla$ 相差一个带基点规范变换（based gauge transformation）$g$。对应于 $\tilde\nabla$ 的流为 $\tilde u_t=g\cdot u_t$，而由于 $g(\tilde p)=I$，有 $\tilde h_t=h_t$。因此 $h_t$ 由 $\rho$ 良定义。于是 $\operatorname{Hom}(\pi,\operatorname{SL}_n(\C))$ 上的流由 $\rho_t=h_t\rho h_t^{-1}$ 定义。该结果表明此流给出到 $\mathrm{hol}(\mathbf{m}^{-1}(0)/\mathcal{G}_E(p))$ 的连续收缩。当 $\rho$ 非半单时，流收敛到半单化（semisimplification）。

这一结果对计算模空间的等变上同调有影响 [2, 12, 43]。特别地，定理 4.3.18 蕴涵
\[
H^*_{\operatorname{SU}_n}\bigl(\mathbf{m}^{-1}(0)/\mathcal{G}_E(p)\bigr)\simeq H^*_{\operatorname{SU}_n}\bigl(\operatorname{Hom}(\pi,\operatorname{SL}_n(\C)\bigr)\bigr).
\]\begin{annote}原文此式右端作 $H^*_{\operatorname{SU}_n}(\operatorname{Hom}(\pi,\operatorname{SL}_n(\C))$，圆括号不配对（缺一个右括号）；译文照排原文。[来源：上下文推断]\end{annote}
注意，由于 $\operatorname{SL}_n(\C)/\operatorname{SU}_n$ 可缩，右端可取关于 $\operatorname{SL}_n(\C)$ 的等变上同调。另一方面，定理 4.2.30 蕴涵
\[
H^*_{\operatorname{SU}_n}\bigl(\mathbf{m}^{-1}(0)/\mathcal{G}_E(p)\bigr)=H^*_{\mathcal{G}_E}\bigl(\mathcal{B}^{\min}_E\bigr)\simeq H^*_{\mathcal{G}_E}\bigl(\mathcal{B}^{ss}_E\bigr).
\]
由此可见，表示空间的等变上同调可以通过研究 $\mathcal{B}_E$ 上 YMH 的等变 Morse 理论（equivariant Morse theory）来计算，这正是 [2] 的精神。这之所以复杂，是因为 $\mathcal{B}_E$ 是奇异的；用这一方法已取得一些进展，见 [13, 67]。

图 4.1 给出了 $\mathcal{C}_E$ 的示意图，其中标出其子空间 $\mathcal{C}^{\mathrm{flat}}_E$ 与 $\mathcal{B}_E$，以及已定义的各个流。

